\PassOptionsToPackage{table,xcdraw}{xcolor}
\documentclass[11pt, a4paper]{article}

\usepackage[T1]{fontenc}
\usepackage[utf8]{inputenc}
\usepackage{tgpagella}
\usepackage[spanish, es-noshorthands]{babel}
\usepackage{amsmath, amssymb}
\usepackage{booktabs}
\usepackage{tabularx}
\usepackage[most]{tcolorbox}
\usepackage[margin=2.5cm]{geometry}
\usepackage{fancyhdr}

\pagestyle{fancy}
\fancyhf{}
\setlength{\headheight}{16pt}
\lhead{\textbf{UCB Sede Tarija} | IMT-221}
\chead{Guía de Redacción: Informes IEEE}
\rhead{Circuitos Electrónicos II}
\lfoot{Documento Estudiantil}
\rfoot{Página \thepage}
\renewcommand{\headrulewidth}{0.4pt}
\definecolor{ucbblue}{RGB}{0, 51, 102}

\begin{document}

\begin{center}
    \Large\textbf{\textcolor{ucbblue}{Guía de Redacción: Informes de Laboratorio (Estructura IEEE)}}[cite: 16]
\end{center}

\begin{tcolorbox}[colback=yellow!5, colframe=ucbblue, title=\textbf{El Formato IEEE en este Curso}]
    La plantilla IEEE es un \textbf{formato de comunicación}, no un sustituto del contenido técnico[cite: 16]. El informe debe demostrar la reconstrucción del siguiente ciclo experimental[cite: 16]:
    \begin{equation*}
        \boxed{ \text{predicción} \rightarrow \text{implementación} \rightarrow \text{medición} \rightarrow \text{comparación} \rightarrow \text{diagnóstico} \rightarrow \text{corrección} \rightarrow \text{validación} } \text{[cite: 16]}
    \end{equation*}
\end{tcolorbox}

\section*{1. Relación entre Laboratorio e Informe[cite: 16]}
\textbf{Para el Lab 3:} El informe debe mostrar como mínimo el circuito implementado, la relación ideal, mediciones DC, discrepancia, diagnóstico, compensación/corrección, comparación before/after e interpretación[cite: 16]. \\
\textbf{Para el Lab 4:} Deberá mostrar configuración, ganancia, datasheets, GBW, Slew Rate, estimación de ancho de banda, $SR_{req}$, simulación, evidencia física, comparación de dispositivos y discusión de la limitación observada[cite: 16].

\section*{2. Resultados y Discusión[cite: 16]}
\textbf{Resultados (Responder: ¿Qué ocurrió?)[cite: 16]} \\
Presente evidencia cruda y tratada. Ejemplo: $V_{o,ideal} = 0.500\,\text{V}$, $V_{o,meas} = 0.522\,\text{V}$, $e = 22\,\text{mV}$[cite: 16].

\textbf{Discusión (Responder: ¿Por qué ocurrió y qué evidencia lo respalda?)[cite: 16]} \\
La discusión debe separar hechos, interpretación, evidencia y conclusiones calificadas[cite: 16]:
\begin{itemize}
    \item \textit{Hecho:} Se midieron $0.522\,\text{V}$[cite: 16].
    \item \textit{Interpretación:} El error de $22\,\text{mV}$ podría contener contribuciones de offset y mismatch[cite: 16].
    \item \textit{Evidencia:} Después de la compensación la salida cambió a $0.505\,\text{V}$[cite: 16].
    \item \textit{Conclusión calificada:} La compensación redujo una componente del error, pero no demuestra que todo el error inicial proviniera de $V_{OS}$[cite: 16].
\end{itemize}

\section*{3. Presentación de Ecuaciones y Figuras[cite: 16]}
\textbf{Ecuaciones:} Cada ecuación importante debe aparecer después de introducirse, usar símbolos definidos, indicar unidades y conectarse con el experimento[cite: 16]. 
\begin{equation*}
    SR_{req} = 2\pi f V_p = 2\pi(20\,000)(5) = 0.628\,\text{V}/\mu\text{s} \text{[cite: 16]}
\end{equation*}
\textbf{Figuras:} Una figura debe responder una pregunta técnica (ej. ¿dónde aparece la distorsión?, ¿qué cambió antes/después?)[cite: 16]. No use imágenes como simple decoración[cite: 16].

\section*{4. Uso de Datasheets[cite: 16]}
No escriba simplemente "el LM741 tiene SR de $0.5\,\text{V}/\mu\text{s}$"[cite: 16]. Registre el fabricante, referencia exacta, valor, si es \textit{typical/min/max}, y las condiciones de prueba[cite: 16]. Distinga siempre entre $\text{valor de datasheet} \neq \text{valor exacto del dispositivo físico}$[cite: 16].

\section*{5. Prácticas a Evitar (Qué NO hacer)[cite: 16]}
\begin{itemize}
    \item Copiar la guía de laboratorio como metodología[cite: 16].
    \item Presentar teoría extensa que no se usa o resultados sin unidades[cite: 16].
    \item Concluir "funcionó correctamente" sin criterios o atribuir toda discrepancia al Op-Amp[cite: 16].
    \item Presentar LTSpice como un experimento físico[cite: 16].
    \item Ocultar mediciones que no coinciden con la teoría. \textit{Una discrepancia bien diagnosticada es más valiosa que una medición "perfecta"[cite: 16].}
\end{itemize}

\end{document}
